\scanpage{104}
% archive sheet number: 103
\deleted{(par le choix d'isomorphismes $\unsure{{}^{12}\mathbb Z_\ell|\mathbb Z_\ell|Z_\ell}(1)\simeq\mathbb Z_\ell$ pour $\forall\,\ell$)
on trouve que \unsure{$u$|$n$|on} définit une représentation
du produit semi-direct de $\mathbb Z$ avec $\mathbb Z$, opérant
sur le sous-groupe $\mathbb Z$ par multiplication par $q$. \deleted{\unsure{Du|Donc}}
Ce groupe est engendré par deux éléments
$f$ (correspondant à $\unsure{\mathrm{frob}_q|\mathrm{frob},q|f\circ b_\rho}$) et $g$, satisfaisant la
relation
\[
fgf^{-1}=g^q \tag*{(6.4)}
\]

Or on a le

\uline{Lemme} 6.5. Soient $F,G$ deux automorphismes d'un
vectoriel \inserted{$E$} de dim. finie sur un corps, \deleted{et} $q$
\deleted{supposant que} un entier $>1$ tel qu'on ait \deleted{\unsure{(6.4)}}
\[
FGF^{-1}=G^q
\]
Alors $G$ est quasi-unipotent, i.e. ses valeurs
propres sont toutes des racines de 1, ou encore :
il existe un entier $N\geq1$ tel que $G^N$ soit
unipotent.

\uline{Dém}. L'ens des valeurs propres de $G$ \deleted{et $FGF^{-1}$}
\deleted{\unsure{étant}} (à valeurs dans une clôture alg. $\overline R$ de $R$)
étant \unsure{le|les} \deleted{\unsure{même, il s'|mêmes, il s'}} $=$ que celui de $FGF^{-1}$,
donc le $=$ que celui de $G^q$, il s'ensuit que
c'est une partie \inserted{\deleted{\unsure{finie}} $\Phi$} de $\overline R^*$ stable par l'application
\[
\lambda\longrightarrow\lambda^q
\]
Comme $\Phi$ est fini, il existe donc pour tt
$\lambda\in\Phi$ deux entiers $N''>N'$ \deleted{tels} $\geq0$ tels que}
